% Brief: whole report, no jargon, avoid acronyms, NO citations, include specific key numbers.
\begin{abstract}
Photographs of moving objects blur, and a shorter exposure trades the blur for darkness and noise.
We built a camera system that avoids this trade-off without modifying the camera: a computer monitor
beside it flashes a designed on/off pattern during one long exposure, which makes the blur
reversible. A fan with letters on its blades, driven by a stepper motor at a known speed, is the
subject. We compare a camera that exposes every row at once with one that exposes row by row,
corrected either once per image or row by row. In simulation, our reconstruction recovers a textured
blade to \psnrGsTex~dB against the stopped fan; correcting row by row matches the all-at-once camera
to within \rowsMaxMeanGapA~dB on average over the image; and the fan's speed is recovered from a single
photograph to within \speedErrTex\%. On the real rig, the monitor showed the pattern to within one
millisecond for \realWithinOneMs~of \realBits~flashes, and with the all-at-once camera at 3 to 8
revolutions per minute the reconstructions score \rigGapMin~to \rigGapMax~dB higher than steady light
reconstructed the same way; the row-by-row camera, which images the fan larger, scores lower but keeps
a \rigRsGapMin~to \rigRsGapMax~dB lead. At 8 revolutions per minute a character reader reads
\rigReadCodedEight\% of the letters after reconstruction with coded light and \rigReadSteadyEight\%
with steady light, and the speed measured from the blur is within 5\% of the set speed on
\rigSpeedWithinCoded~of \rigSpeedAllCoded~and \rigSpeedWithinRsCoded~of \rigSpeedAllRsCoded~coded
photographs from the two cameras.
\end{abstract}
